\section{Introduction}
\label{sec:introduction}

Pengenalan otomatis atribut demografis wajah memegang peranan penting dalam berbagai domain aplikasi cerdas, termasuk sistem forensik digital, kontrol akses biometrik, interaksi manusia-komputer, dan personalisasi layanan interaktif. Namun, sejumlah studi melaporkan bahwa sistem pengenalan wajah kerap memperlihatkan disparitas performa terhadap subkelompok demografis tertentu, dengan penurunan akurasi yang substansial secara konsisten teramati pada kelompok wanita dan populasi berkulit gelap. Ketimpangan akurasi tersebut bersumber dari beragam faktor fundamental, meliputi ketidakseimbangan distribusi data latih pada dataset publik dan keterbatasan ekstraktor fitur konvensional dalam mengekstraksi representasi visual yang invarian terhadap variasi pose serta degradasi visual alami. Keterbatasan sistem konvensional dalam mengatasi variasi visual ini menegaskan bahwa pengenalan demografis yang adil dan akurat masih menghadapi kendala besar, sehingga diperlukan pendekatan baru yang mampu memodelkan atribut demografis secara komprehensif.

Sebagian besar literatur pengenalan atribut wajah memformulasikan klasifikasi ras dan gender sebagai dua tugas terpisah yang diselesaikan secara independen. Pendekatan terisolasi ini mengabaikan sifat alami identitas visual wajah dan tidak mengevaluasi performa pengenalan secara eksplisit pada persilangan kelompok demografis, seperti wanita berkulit hitam atau wanita Asia. Akibatnya, penurunan performa yang terkonsentrasi pada subkelompok interseksional tertentu luput dari pengamatan metrik global, sehingga kondisi tersebut membatasi pemahaman komprehensif mengenai keandalan model di lapangan. Formulasi klasifikasi interseksional terpadu pada pengaturan evaluasi seimbang menyediakan landasan komparatif yang objektif untuk mengukur variasi performa lintas subkelompok secara adil. Namun, pemodelan klasifikasi interseksional menuntut representasi visual yang kaya dan berkapasitas tinggi untuk memisahkan batas keputusan antarsubkelompok yang memiliki derajat tumpang tindih visual tinggi, sehingga pemodelan ini memerlukan arsitektur ekstraksi representasi visual yang tangguh terhadap variasi fenotipe dan ekspresi wajah.

Perkembangan metodologi pengenalan atribut wajah telah bertransformasi secara substansial dari pendekatan ekstraksi fitur buatan tangan menuju arsitektur deep learning dan ViT. Pada era arsitektur Convolutional Neural Networks (CNN), berbagai penelitian mengeksplorasi ekstraksi ciri demografis melalui pemanfaatan region spesifik wajah seperti area wajah bagian tengah \cite{ref1}, kerangka kerja multi-tugas berbasis multi-skala untuk estimasi usia dan gender secara simultan \cite{ref2}, serta model CNN mendalam untuk pengenalan etnisitas \cite{ref3}. Perkembangan terkini menunjukkan keunggulan ViT dalam memodelkan ketergantungan spasial global antarpatch citra melalui mekanisme Multi-Head Self-Attention (MHSA), yang membuka ruang analisis mendalam terhadap bias representasi visual \cite{ref4}, sementara pemanfaatan representasi generatif sintetis turut dikembangkan untuk memitigasi disparitas demografis \cite{ref5}. Di samping itu, integrasi mekanisme atensi paralel \cite{ref6}, pemanfaatan ViT untuk klasifikasi gender wajah \cite{ref7}, serta eksplorasi arsitektur transformer multi-aksis \cite{ref8} semakin memperkaya variasi pemodelan representasi visual. Lebih lanjut, eksplorasi penggabungan representasi ViT lintas domain mulai dikembangkan melalui integrasi fitur ganda wajah dan emosi untuk meningkatkan diskriminasi atribut ras dan gender \cite{ref9}, \cite{ref10}. Meskipun demikian, eksplorasi fusi representasi multi-domain komprehensif yang secara simultan mengintegrasikan isyarat biometrik struktural, dinamika ekspresi emosi, dan morfologi penuaan biologis wajah masih memerlukan kajian empiris terpadu.

Meskipun eksplorasi ViT dan model visi telah berkembang pesat, terdapat tiga kesenjangan penelitian utama yang belum terpecahkan secara optimal. Pertama, sebagian besar studi terdahulu bertumpu pada representasi biometrik wajah domain tunggal tanpa mengintegrasikan isyarat afektif dinamis secara simultan, sehingga model yang dikembangkan kesulitan membedakan subkelompok dengan derajat phenotypic overlap yang tinggi. Kedua, pengenalan demografis sering kali menunjukkan sensitivitas tinggi terhadap perubahan morfologi wajah seiring pertambahan usia karena ekstraksi fitur yang digunakan tidak mengikutsertakan representasi penuaan biologis. Ketiga, evaluasi komparatif mengenai perilaku batas keputusan antarmodel classifier hilir yang mencakup model linier, probabilistik, ensemble, dan kernel-based, beserta interaksinya terhadap penskalaan dan reduksi dimensi pada ruang representasi transformer, masih sangat terbatas. Untuk menjembatani keterbatasan representasi domain tunggal serta mengevaluasi batas keputusan classifier hilir tersebut, penelitian ini mengusulkan integrasi tiga domain representasi ViT yang dipadukan dengan optimasi pipeline classical classifier terstruktur.

Untuk mengatasi kesenjangan riset tersebut, penelitian ini mengusulkan kerangka kerja fusi fitur Tri-Domain ViT yang memadukan tiga representasi laten spesifik tugas untuk klasifikasi demografis interseksional pada citra wajah. Kerangka kerja yang diusulkan mengintegrasikan tiga domain representasi komplementer, meliputi representasi geometri biometrik wajah, isyarat ekspresi afektif dinamis, dan karakteristik morfologi penuaan biologis dari model backbone ViT pre-trained spesifik tugas ke dalam satu representasi terpadu melalui operasi konkatenasi fitur. Ekstraksi fitur laten dilakukan secara offline dari model backbone pre-trained yang dibekukan, sehingga prosedur ini mampu menjaga efisiensi komputasi sekaligus mengisolasi representasi visual dari variabilitas proses pelatihan ulang. Vektor fitur yang dihasilkan selanjutnya dievaluasi secara sistematis melalui serangkaian skema ablasi domain tunggal, domain ganda, dan tri-domain pada empat kelompok algoritma classical machine learning, yaitu RF, GNB, LR, dan SVM. Seluruh arsitektur klasifikasi hilir dioptimalkan melalui prosedur Stratified Cross-Validation yang dikombinasikan dengan pipeline modular penskalaan fitur dan reduksi dimensionalitas Principal Component Analysis (PCA) terisolasi, untuk memastikan pencegahan kebocoran data dan pembentukan batas keputusan yang tangguh antarsubkelompok demografis.

Kontribusi ilmiah utama dari penelitian ini dirangkum ke dalam empat poin utama sebagai berikut:
\begin{enumerate}
\item A Tri-Domain ViT feature fusion framework integrating face-associated biometric representations, expression-related representations, and age-associated facial representations into a unified latent feature vector for intersectional demographic classification (Kerangka kerja fusi fitur Tri-Domain ViT yang mengintegrasikan representasi terkait biometrik wajah, representasi terkait ekspresi, dan representasi terkait usia wajah ke dalam vektor fitur laten terpadu untuk klasifikasi demografis interseksional).
\item An empirical comparative benchmark across multiple feature ablation schemes and classical machine learning classifiers optimized via Stratified Cross-Validation hyperparameter tuning, examining decision boundary behavior across linear, probabilistic, ensemble, and kernel-based models (Tolok ukur komparatif empiris melintasi berbagai skema ablasi fitur dan pengklasifikasi machine learning klasik yang dioptimalkan melalui penyetelan hyperparameter validasi silang berstrata, mengkaji perilaku batas keputusan antarmodel linier, probabilistik, ensemble, dan berbasis kernel).
\item Competitive classification performance on the DemogPairs dataset, achieving higher reported performance compared to the evaluated single-domain and dual-domain configurations on the majority of classifiers (Capaian performa klasifikasi kompetitif pada dataset DemogPairs yang mencapai performa lebih tinggi dibandingkan konfigurasi domain tunggal dan domain ganda yang dievaluasi pada mayoritas pengklasifikasi).
\item A granular intersectional subgroup performance analysis, evaluating subgroup-level classification performance across six demographic subgroups, providing an in-depth characterization of model behaviors across racial and gender intersections while systematically revealing the consistency of multi-domain feature representations on independent evaluation cohorts without exhibiting acute performance degradation (Analisis performa subkelompok interseksional granular yang mengevaluasi performa klasifikasi di keenam subkelompok demografis, memberikan karakterisasi mendalam mengenai perilaku model melintasi persilangan ras dan gender sekaligus mengungkap secara sistematis konsistensi representasi fitur multi-domain pada kohort evaluasi independen tanpa memperlihatkan degradasi performa yang akut).
\end{enumerate}

Sistematika penulisan artikel ini disusun sebagai berikut: Section~\ref{sec:related_works} mengulas tinjauan pustaka mengenai pengenalan demografis wajah dan arsitektur transformer. Section~\ref{sec:materials_and_methods} menjabarkan spesifikasi dataset DemogPairs, ekstraksi representasi multi-domain ViT, dan pipeline klasifikasi teroptimasi. Section~\ref{sec:results_and_discussion} memaparkan hasil evaluasi eksperimental, studi ablasi fitur, serta analisis performa subkelompok interseksional. Terakhir, Section~\ref{sec:conclusion} merangkum kesimpulan utama penelitian, keterbatasan metodologis, dan rekomendasi arah riset mendatang.
